\section{Related Work}

\paragraph{Visual SLAM and Dense Pose Estimation.}
Visual SLAM systems estimate camera motion while simultaneously reconstructing scene geometry through feature-based or dense optimization pipelines \cite{murartal2015orbslam, engel2014lsdslam, klein2007ptam, newcombe2011kinectfusion}. These approaches typically rely on persistent mapping, global bundle adjustment, and long-term state estimation. While highly accurate, such pipelines incur substantial computational and memory overhead, particularly for dense direct methods that operate on per-pixel photometric alignment. As a result, deploying full SLAM systems on mobile or embedded devices often introduces significant latency and power consumption, making them unsuitable for real-time augmented reality scenarios where only object-level pose tracking is required.

\paragraph{Model-Based Tracking for Textureless Objects.}
Model-based approaches leverage known object geometry to estimate pose by aligning projected model features with image observations. A representative example is the edge-based tracking framework proposed by Choi and Christensen, which aligns projected contours with image gradients using a particle filtering formulation \cite{choi2012textureless}. Such methods are effective for textureless objects and exploit strong geometric priors. However, they typically rely on dense edge measurements and iterative evaluation of multiple pose hypotheses, resulting in high computational cost. More broadly, edge- and contour-based tracking methods remain sensitive to noise, clutter, and partial occlusion, limiting their robustness in real-world mobile scenarios \cite{han2019review}.

\paragraph{Feature-Based Object Tracking.}
Classical feature-based pipelines estimate motion by detecting keypoints, computing local descriptors, and matching them across frames. Efficient binary descriptors such as ORB \cite{rublee2011orb} and BRISK \cite{leutenegger2011brisk} enable fast matching using Hamming distance and are widely used in real-time systems \cite{rosten2006fast, lowe2004sift}. Pose is typically recovered using algebraic formulations such as planar homography estimation or Perspective-$n$-Point (PnP) solvers \cite{kneip2011pnp, gao2003pnp, hartley2003multiple}. While computationally efficient, these approaches rely purely on image-space correspondences and lack explicit geometric reasoning. Consequently, they often exhibit instability under large viewpoint changes, rapid scale variations, and planar degeneracies.

\paragraph{Learning-Based Pose Estimation.}
Recent work has explored deep neural networks that regress object pose directly from RGB images. These methods can achieve robustness to appearance variation, occlusion, and illumination changes when trained on large annotated datasets. However, they require substantial training data and computational resources, and inference often depends on GPU acceleration. Furthermore, adapting such models to new objects typically requires retraining or fine-tuning, limiting their flexibility in deployment scenarios with constrained compute budgets.

\paragraph{Render-and-Compare Pose Tracking.}
Render-and-compare approaches estimate pose by rendering synthetic views of known object geometry and minimizing discrepancies with the observed image. By incorporating explicit geometric reasoning, these methods can achieve stable pose estimation through photometric or geometric alignment. However, dense render-and-optimize formulations typically involve iterative optimization over high-dimensional image representations, making them computationally expensive and difficult to deploy efficiently on mobile hardware.

\paragraph{Optimization-Based Pose Refinement.}
Pose refinement is commonly formulated as a nonlinear least-squares problem on the $\mathrm{SE}(3)$ manifold and solved using Gauss--Newton or Levenberg--Marquardt optimization \cite{triggs1999bundle, nocedal2006}. While effective, these methods treat all correspondences uniformly or rely on predefined robust loss functions (e.g., Huber or Cauchy) to mitigate outliers \cite{huber1964, black1996}. Such approaches do not explicitly exploit the relative structure of residuals within a given frame, and often require careful tuning of loss parameters to balance robustness and convergence.

\paragraph{Hybrid Tracking Frameworks.}
To balance robustness and efficiency, hybrid systems combine global feature matching with lightweight local tracking methods such as Lucas--Kanade optical flow. These approaches enable reliable relocalization while maintaining efficient frame-to-frame tracking. Building on this principle, our work explores a hierarchy of tracking configurations tailored for edge devices: (i) a SIMD-optimized CPU pipeline for efficient feature processing, (ii) a GPU-accelerated pipeline using programmable shaders for feature extraction and parallel matching, and (iii) a geometry-aware $\mathrm{SE}(3)$ render-and-track framework that integrates synthetic reference rendering with adaptive, correspondence-driven pose refinement.